\chapter{總結}

本論文主要探討如何利用金屬邊緣電漿子（Edge Plasmons, EP）實現動量工程，藉此克服傳統光學激發受限於光錐範圍的瓶頸，以耦合單層過渡金屬二硫族化合物（TMDs）中的高動量谷激子。我們利用時域有限差分法（FDTD）與數值分析，探討了金屬薄膜厚度對 EP 色散關係與近場分佈的影響，並計算其與 TMD 激子的近場耦合矩陣元。

在色散關係的分析中，模擬結果顯示 Au 薄膜的厚度會直接改變 EP 的色散行為。隨著薄膜厚度由 40\,nm 減小至 10\,nm，邊緣電荷間的靜電恢復力減弱，導致 EP 的本徵頻率發生紅移，色散曲線也隨之變得較為平坦。這種幾何調控效應使 EP 在相同的激發能量下，能夠提供比傳統表面電漿極化子（SPP）或自由空間光子更大的波向量，具備較佳的動量補償能力。

在近場耦合方面，我們擷取了動量空間中的 EP 向量勢 $\mathbf{A}(\mathbf{Q})$，並代入包含偽自旋（Pseudo-spin）選擇定則的近場耦合矩陣元 $\tilde{M}_{S,\mathbf{Q}}$ 中進行計算。結果表明，由於空間偏振分量的相消干涉，$K$ 谷激子與 EP 的耦合強度在動量空間中心軸（$Q_x = 0$）處會出現明顯的低耦合強度暗帶。這顯示出 TMD 激子的谷相依特性會直接影響耦合分佈，且最大耦合強度會發生在偏離中心的橫向動量 $Q_x^*$ 處。

進一步將 EP 的色散關係與 $\text{MoSe}_2$ 的激子色散曲線疊加後，可發現兩者的共振交點會隨著金屬厚度的改變而移動。當厚度減薄時，滿足能量與動量守恆的交會點 $(E^*, Q^*)$ 會逐漸往高動量區域偏移。這項結果證實了改變單一的幾何參數（厚度）即可主動調控共振條件，實現微觀尺度下的動量工程。

綜合上述結果，本論文建立了邊緣電漿子與 TMD 激子耦合的數值與理論分析模型，並證實了利用金屬厚度進行動量調控的可行性。這項研究結果有助於理解超越光錐極限的二維材料光學響應，並可作為未來開發基於激子-電漿子強耦合的室溫光電元件或谷微納光子學應用時的結構設計參考。